\section{总结与延伸}

\subsection{一条贯穿全课的主线：控制模型能看见的捷径}

两篇工作看似一个讲 object，一个讲 Gaussian regularization，实则都在重新设计可行解空间。Causal-JEPA 删除目标对象自身历史，让 interaction shortcut 之外的关系信息变得必要；LeWorldModel 约束 embedding distribution，让常数表示不再是 prediction loss 的合法最优解。高质量 self-supervised learning 的共同原则，就是让``便宜但错误''的解难以满足目标。

\begin{importantbox}{本课最可迁移的八问}
\begin{enumerate}
  \item 状态表示保留了哪些决策相关变量？
  \item 最容易利用的 shortcut 是什么？
  \item anti-collapse 约束排除了哪些平凡解？
  \item 动作以什么接口进入 predictor？
  \item 训练时可见的信息与推理时是否一致？
  \item 结果增益来自表示、architecture、loss 还是 controller？
  \item probe、attention 与漂亮 rollout 各自不能证明什么？
  \item 长 horizon、分布外动作和真实环境怎样失败？
\end{enumerate}
\end{importantbox}

\subsection{结构偏置与优化偏置必须协同}

只有 object masking、没有稳定 representation，模型可能在退化 latent 上学习所谓 interaction；只有 SIGReg、没有任务结构，模型可能保持丰富方差却仍依赖背景相关性。下一代 world model 需要同时设计 representation units、observability interventions、anti-collapse distribution、action interface 与 planner。它们应在同一 evidence protocol 中被逐项 ablate。

\begin{knowledgebox}{Causal-JEPA 与 LeWorldModel 的互补关系}
Causal-JEPA 主要回答 ``预测应依赖哪些变量''；LeWorldModel 主要回答 ``encoder 怎样保持有信息''。前者目前依赖冻结 object-centric encoder，后者目前使用全局 compact latent。一个自然延伸是端到端学习稳定 object slots，同时保留 interaction-forcing masks 与可扩展 anti-collapse regularizer。
\end{knowledgebox}

\subsection{三条证据边界}

第一，controlled benchmarks 中的 counterfactual、planning 与 surprise gains 不等于开放世界 causal understanding。第二，latent 可解码与 attention 集中不等于 predictor 真正使用了正确机制。第三，课堂时点的论文标题和理论表述会随审稿与后续分析修订；研究讲义应保存历史快照，也应标出后来更谨慎的边界。

\begin{warningbox}{不要把世界模型写成万能中间层}
世界模型会压缩观测，压缩就会丢信息；会 rollout，rollout 就会累积误差；会被 planner 调用，planner 就可能 exploit model error。它是把预测、控制与验证连接起来的强接口，不是自动消除 hallucination、safety 与 distribution shift 的魔法层。
\end{warningbox}

\subsection{开放问题}

值得继续追踪的方向包括：如何在真实视频中稳定发现和追踪对象；如何把多个时间尺度与事件抽象放进 latent；如何用 uncertainty 标记模型不知道的未来；如何让 language goal、reward、constraint 与 image goal 共存；如何让 planner 不利用 model error；以及如何用真实机器人干预区分相关性、predictive sufficiency 与可识别 causal mechanism。

\subsection{拓展阅读}

建议按课堂时点阅读 Causal-JEPA v1 与 LeWorldModel v1，再对照 Causal-JEPA 后续修订，观察 causal wording 如何变得更谨慎。背景材料可依次阅读 Slot Attention、V-JEPA、V-JEPA 2 与 DINO-WM；实现层面可查看两篇工作的官方代码和 `stable-worldmodel` 统一实验库。阅读时继续使用本课的八问，而不是只比较 leaderboard 数字。

\end{document}
